\section*{Bab \ref{ch:b1:induction} --- Induksi Elektromagnetik dan Penerapannya}

\begin{solution}{exo:b1:induction:1}
$\Phi = NBS\cos30^\circ = 100 \times 0.2 \times 10^{-3} \times 0.866 = \qty{1.7e-2}{Wb}$;
$\langle e\rangle = \Delta\Phi/\Delta t = \qty{1.7}{V}$ dalam \qty{10}{ms}, \qty{17}{V} dalam
\qty{1}{ms}. Arusnya mengalir sedemikian sehingga menciptakan kembali fluks yang lenyap:
medannya menyusuri $\vect B$ yang lama.
\end{solution}

\begin{solution}{exo:b1:induction:2}
$e = B\ell v = 0.5 \times 0.2 \times 2 = \qty{0.20}{V}$; $i = \qty{2.0}{A}$; $F = B\ell i =
\qty{0.20}{N}$ yang mengerem, sehingga pendorongnya mengerahkan \qty{0.20}{N}; $Fv = \qty{0.40}{W}
= Ri^2$.
\end{solution}

\begin{solution}{exo:b1:induction:3}
$L = \mu_0n^2S\ell = 4\pi \times 10^{-7} \times 10^6 \times 4 \times 10^{-4} \times 0.2 =
\qty{0.10}{mH}$; $W = \tfrac12LI^2 = \qty{1.3}{mJ}$; $\tau = L/R = \qty{0.1}{ms}$.
$B = \mu_0nI = \qty{6.3}{mT}$, $B^2/2\mu_0 = \qty{15.7}{J/m^3}$ atas
$S\ell = \qty{8e-5}{m^3}$: \qty{1.3}{mJ}.
\end{solution}

\begin{solution}{exo:b1:induction:4}
$N_2/N_1 = 12/230 = 0.052$; $N_1 = 50/0.052 = 960$ lilitan; $i_1 = 2 \times 0.052 =
\qty{0.10}{A}$. Arus tunak membuat fluks tunak: tak ada $\dd\varphi/\dd t$, tak ada
\omterm{def:b1:dc-circuits:voltage}{tegangan} sekunder --- dan primernya, sebagai hubung singkat, akan terbakar.
\end{solution}

\begin{solution}{exo:b1:induction:5}
(a) Fluksnya (menyusuri arah datangnya) tumbuh: arusnya membuat medan
yang melawannya --- kutub utaranya sendiri menghadapi kutub utara magnetnya --- dan
cincinnya terdorong menjauh. (b) Arus yang medannya melawan pertumbuhannya;
cincinnya, sebagai dipol yang berlawanan arah dengan $\vect B$, terdorong ke medan
yang lebih lemah. (c) Selagi jatuh fluksnya tumbuh: arusnya sama seperti (a), cincinnya
tererem. (d) Tiap penggal tabung di depan magnetnya melihat fluks yang tumbuh,
di belakangnya fluks yang mengecil; kedua arus imbasnya mengerem magnetnya; pada
kelajuan yang membuat \omterm{thm:b1:magnetostatics:laplace}{gaya Laplace} mengimbangi beratnya, magnetnya jatuh
tunak dan \omterm{def:b1:work-and-energy:conservative}{energi potensial} yang hilang menjadi kalor Joule dalam tabungnya.
\end{solution}

\begin{solution}{exo:b1:induction:6}
$E_{\max} = NBS\omega = 200 \times 2 \times 10^{-3} \times 0.1 \times 314 = \qty{12.6}{V}$, rms-nya
\qty{8.9}{V}; $I_{\max} = \qty{1.26}{A}$, $\Gamma_{\max} = NBSI_{\max} = \qty{0.050}{N.m}$
(ketika $\sin\omega t = 1$); $\langle P\rangle = E_{\max}^2/2R = \qty{7.9}{W}$ $=
\langle\Gamma\omega\rangle = \tfrac12 \times 0.05 \times 314$; ketika $e = 0$ arusnya
nol dan \omterm{def:b1:angular-momentum:moment}{momen gayanya} lenyap --- \omterm{def:b1:angular-momentum:moment}{momen gayanya} berdenyut pada \qty{100}{Hz}.
\end{solution}

\begin{solution}{exo:b1:induction:7}
$e = Bav$, $i = Bav/R$, $F = B^2a^2v/R = 10\,v$ (N). $m\,\dd v/\dd t = -B^2a^2v/R$:
$v = v_0\eu^{-t/\tau}$, $\tau = mR/B^2a^2 = \qty{0.10}{s}$; jaraknya $v_0\tau =
\qty{10}{cm}$ (bila daerah medannya cukup panjang). \omterm{thm:b1:work-and-energy:ke}{Energi kinetiknya},
\qty{0.5}{J}, dilesapkan sebagai $Ri^2$ dalam lilitannya.
\end{solution}

\begin{solution}{exo:b1:induction:8}
Fluks medan solenoidanya $\mu_0n_1i_1$ menembus $N_2$ lilitannya: $M =
\mu_0n_1N_2S_2 = 4\pi \times 10^{-7} \times 2000 \times 50 \times 2 \times 10^{-4} = \qty{25}{\micro H}$.
Tanjakannya: $e = M\,\dd i_1/\dd t = \qty{2.5}{mV}$; pada \qty{10}{kHz}, \qty{1}{A}: $e_{\max} =
M\omega I = 25 \times 10^{-6} \times 6.3 \times 10^4 = \qty{1.6}{V}$. $M$ bersifat setangkup:
angka yang sama memberi fluks yang dikirim kumparan kecilnya menembus
solenoidanya --- yang jauh lebih sukar dihitung langsung.
\end{solution}

\begin{solution}{exo:b1:induction:9}
$W = B^2V/2\mu_0 = 2.25/(2.51 \times 10^{-6}) = \qty{0.90}{MJ}$; $L = 2W/I^2 =
\qty{7.2}{H}$; $h = W/mg = \qty{92}{m}$; heliumnya: $9 \times 10^5/2600 = \qty{350}{L}$.
Medan Bumi: $(5 \times 10^{-5})^2/2\mu_0 = \qty{1e-3}{J/m^3}$, \qty{50}{mJ} dalam
ruangannya.
\end{solution}

\begin{solution}{exo:b1:induction:10}
(a) $E - B\ell v = Ri$ (ggl baliknya melawan sumbernya), $m\,\dd v/\dd t =
B\ell i$: $m\,\dd v/\dd t = B\ell(E - B\ell v)/R$. (b) $v_\infty = E/B\ell$, $\tau =
mR/B^2\ell^2$. (c) Kalikanlah persamaan listriknya dengan $i$: $Ei = Ri^2 +
B\ell vi$, dan $B\ell vi = Fv$ adalah daya mekanisnya; \omterm{def:b1:heat-engines:efficiency}{efisiensinya} $Fv/Ei =
B\ell v/E = v/v_\infty$ --- $100\%$ hanya pada keadaan tanpa beban, yang tak
menyalurkan apa pun. (d) $v_\infty = 2/(0.5 \times 0.2) = \qty{20}{m/s}$, $\tau = 0.05 \times
0.1/0.01 = \qty{0.5}{s}$. (e) $F = B\ell I = 1 \times 0.1 \times 10^6 = \qty{1e5}{N}$
(dengan mengambil $\ell = \qty{10}{cm}$), $a = \qty{1e7}{m/s^2}$, $v = \sqrt{2a \times 5} =
\qty{1e4}{m/s}$ --- sepuluh kilometer tiap sekon, janji sekaligus
masalah senapan rel (relnya terkikis).
\end{solution}

\begin{solution}{exo:b1:induction:11}
(a) Pada jari-jari $r$ logamnya bergerak dengan $v = \omega r$; $\vect v\wedge\vect B$
radial dan bernilai $\omega rB$: $e = \int_0^R\omega Br\,\dd r = \tfrac12B\omega R^2$.
(b) $\omega = 314$: $e = 0.5 \times 1 \times 314 \times 0.01 = \qty{1.6}{V}$. (c) Hanya satu
``lilitan'', jadi paling banter beberapa volt; ke dalam \qty{1}{m\ohm}: \qty{1.6}{kA}, \omterm{def:b1:angular-momentum:moment}{momen gayanya}
$= P/\omega = ei/\omega = 2500/314 = \qty{8}{N.m}$. (d) Rangkaian sikat--jari-jari--
tepi--sikatnya menyapu luas $\tfrac12R^2\omega\,\dd t$ tiap $\dd t$: $\dd\Phi/\dd t =
\tfrac12BR^2\omega$ --- yaitu fluks menembus rangkaian yang bergerak dan berubah bentuk itu.
\end{solution}

\begin{solution}{exo:b1:induction:12}
(a) $e_2 = -M\,\dd i_1/\dd t$, puncaknya $M\omega I_1 = 10^{-6} \times 1.57 \times 10^5 \times 30 =
\qty{4.7}{V}$. (b) $I_2 = 4.7/0.01 = \qty{470}{A}$ puncak; $\langle P\rangle = R_pI_2^2/2
= \qty{1.1}{kW}$. (c) Gglnya $\propto\omega$: frekuensi tinggi memberi ggl
yang berguna dengan $M$ yang sederhana; kacanya isolator, tak ada yang mengalir di dalamnya
--- kalornya lahir di dalam logamnya. (d) Arus pancinya mengimbas pada
kumparannya $M\omega I_2 = 10^{-6} \times 1.57 \times 10^5 \times 470 = \qty{74}{V}$ puncak,
sefase dengan $i_1$ (arus pancinya tak tertinggal terhadap $e_2$, karena
bersifat resistif, dan $e_2$ tertinggal terhadap $i_1$ sebesar \qty{90}{\degree}\ldots\ sehingga
ggl pantulnya $-M\,\dd i_2/\dd t$, berlawanan fase dengan \omterm{def:b1:dc-circuits:voltage}{tegangan}
\emph{penggerak} $i_1$): generatornya memasok $\tfrac12 \times 74 \times 30 =
\qty{1.1}{kW}$ lebih banyak --- persis yang dilesapkan pancinya.
\end{solution}

\begin{solution}{pb:b1:induction:1}
\textbf{1.} $\Phi = NBS\cos\omega t$, $e = NBS\omega\sin\omega t$; $NBS = 300 \times
0.3 \times 3 \times 10^{-4} = \qty{0.027}{Wb}$; $\omega = v/r = 5/0.01 = \qty{500}{rad/s}$:
\omterm{def:b1:wave-propagation:signal}{amplitudonya} \qty{13.5}{V}, frekuensinya $\omega/2\pi = \qty{80}{Hz}$.

\textbf{2.} $I_{\text{rms}} = 13.5/\sqrt2/16 = \qty{0.60}{A}$; lampunya $RI^2 = \qty{4.3}{W}$
di bawah \qty{7.2}{V}: lampu \qty{6}{V}, \qty{3}{W} terpacu berlebihan dan tak akan
bertahan.

\textbf{3.} Mekanis $=$ listrik: $(R + r_c)I_{\text{rms}}^2 = 16 \times 0.36 =
\qty{5.7}{W}$; $F = P/v = 5.7/5 = \qty{1.1}{N}$ --- separuh gesekan gelindingnya;
pengendaranya merasakannya.

\textbf{4.} Kopel Laplace arus imbasnya melawan putarannya
(Lenz); energi lampunya datang dari kaki pengendaranya, sebanyak \qty{5.7}{W}.

\textbf{5.} $\underline Z = R + r_c + jL\omega$; $I = NBS\omega/\sqrt{(R + r_c)^2 + L^2\omega^2}$;
untuk $L\omega \gg R + r_c$, $I \to NBS/L = 0.027/0.02 = \qty{1.35}{A}$. $70\%$: $L\omega
= 0.7(R + r_c)/\sqrt{1 - 0.49} = 15.7$, $\omega = \qty{784}{rad/s}$, $v = \qty{7.8}{m/s}
= \qty{28}{km/h}$.

\textbf{6.} Tanpa $L$: $e_{\max} = 3 \times 13.5 = \qty{40.5}{V}$, $I_{\text{rms}} =
\qty{1.8}{A}$, \qty{39}{W} pada lampunya --- ia putus. Dengan $L$: $I = 40.5/\sqrt{256 +
900} = \qty{1.2}{A}$ puncak, \qty{8.5}{W}: panas, tetapi sepuluh kali lebih kecil.

\textbf{7.} \qty{9}{km/h}: $e_{\max} = 6.75$, $|Z| = \sqrt{256 + 25} = 16.8$, $I = 0.40$,
$P = \tfrac12 \times 0.16 \times 12 = \qty{1.0}{W}$: redup. \qty{36}{km/h}: $e_{\max} = 27$,
$|Z| = 25.6$, $I = 1.05$, $P = \qty{6.6}{W}$. Kelajuannya empat kali lipat, arusnya $2.6$ kali:
induktansinya mengatur.

\textbf{8.} Hanya gerak nisbinya yang berperan: fluks menembus kumparan
yang tetap berubah dengan cara yang sama, gglnya sama. Tak ada sentuhan luncur (sikat) yang
diperlukan untuk membawa arusnya keluar dari kumparan yang berputar.

\textbf{9.} Tak ada gerak, tak ada perubahan fluks, tak ada cahaya --- karena itu ada \omterm{def:b1:potential-capacitors:capacitance}{kapasitor} atau
baterai kecil (``lampu diam'') yang diisi selagi bersepeda.

\textbf{10.} Tiap unsur kawatnya tegak lurus pada $\vect B$ yang radial:
$F = B\ell i$, aksial. Bergerak dengan $v$, tiap unsurnya melihat $\vect v\wedge\vect B$
menyusuri kawatnya: $e = -B\ell v$, yang melawan arus penghasil
geraknya (\omterm{thm:b1:dc-circuits:kirchhoff}{kesepakatan penerima}: \omterm{def:b1:dc-circuits:voltage}{tegangan} $B\ell v$).

\textbf{11.} $u = Ri + L\,\dd i/\dd t + B\ell v$; $m\,\dd v/\dd t = -kx - hv + B\ell i$.

\textbf{12.} $\underline i = (\underline u - B\ell\underline v)/(R + jL\omega)$; $\underline v\,(jm\omega + h +
k/j\omega) = B\ell\,\underline i$; substitusikan lalu selesaikan untuk $\underline v$.

\textbf{13.} $i \approx (u - B\ell v)/R$ memberi $m\dot v + (h + B^2\ell^2/R)v + kx =
B\ell u/R$: peredaman tambahannya $25/8 = \qty{3.1}{N.s/m}$. $f_0 = \tfrac1{2\pi}\sqrt{k/m} =
\qty{50}{Hz}$; $Q = \sqrt{km}/h_{\text{tot}}$: $3.2$ dalam keadaan terbuka, $0.77$ dengan
penguatnya.

\textbf{14.} $F = \qty{5.0}{N}$, $a = F/m = \qty{500}{m/s^2}$, $x = a/\omega^2 =
\qty{13}{\micro m}$, $v = a/\omega = \qty{0.080}{m/s}$; $B\ell v = \qty{0.40}{V}$ terhadap
$Ri = \qty{8}{V}$: $5\%$.

\textbf{15.} Dengan arusnya dipaksakan, $v = F/h = \qty{5}{m/s}$ (yaitu
\omterm{def:b1:wave-propagation:signal}{amplitudo} \qty{16}{mm} --- simpangan yang besar untuk \qty{1}{A}); $B\ell v =
\qty{25}{V}$, tiga kali $Ri$: penguatnya harus memasok \qty{33}{V} untuk mendorong
\qty{1}{A} pada resonansinya --- itulah puncak impedansnya. Sebuah penguat bersumber \omterm{def:b1:dc-circuits:voltage}{tegangan}
justru akan membiarkan peredaman listriknya menahan kerucutnya.

\textbf{16.} $\tfrac12RI^2 = \qty{4.0}{W}$; $\tfrac12hv^2 = \tfrac12 \times 0.0064 =
\qty{3.2}{mW}$: \omterm{def:b1:heat-engines:efficiency}{efisiensinya} $0.08\%$ --- \omterm{prop:b1:induction:loudspeaker}{pengeras suara} adalah pemanas yang
berbisik.

\textbf{17.} Di atas resonansinya $a = F/m = B\ell i/m$ tak bergantung pada
frekuensi bagi arus tertentu, sehingga tekanannya, yang $\propto a$, mendatar:
pita yang terkendali massa itulah pita kerjanya.

\textbf{18.} Medan radial: gayanya aksial dan sama di setiap
kedudukan kumparannya di dalam celahnya. Sebuah tweeter harus tetap terkendali massa
sampai \qty{20}{kHz} dengan simpangan yang mungil: kumparan dan kerucut yang ringan; dan
kerucut yang kecil memancarkan frekuensi tinggi ke segala arah.

\textbf{19.} Frekuensi rendah: $v$ kecil, $|Z| \approx R = \qty{8}{\ohm}$; pada
\qty{50}{Hz}: $u = Ri + B\ell v$ dengan $v = B\ell i/h$: $Z = R + B^2\ell^2/h = 8 + 25 =
\qty{33}{\ohm}$; pada \qty{10}{kHz}: $|R + jL\omega| = |8 + 31j| = \qty{32}{\ohm}$. Sebuah
puncak pada resonansinya, lantai datar \qty{8}{\ohm}, lalu naik bersama $L\omega$.

\textbf{20.} $e = B\ell v = 5 \times 10^{-3} = \qty{5}{mV}$.

\textbf{21.} $i = 5 \times 10^{-3}/608 = \qty{8.2}{\micro A}$; $F = B\ell i = \qty{41}{\micro N}$,
yang melawan $v$: $F = [B^2\ell^2/(R + R_L)]\,v$, yaitu peredaman \qty{0.041}{N.s/m};
daya $Fv = \qty{4e-8}{W}$ diambil dari bunyinya dan berakhir di
kedua \omterm{def:b1:dc-circuits:resistor}{resistornya}.

\textbf{22.} Kelosnya: $B^2\ell_f^2/R_f = 1 \times 0.01/10^{-3} = \qty{10}{N.s/m}$, sepuluh
kali $h$: ia akan mematikan kerucutnya dan memboroskan daya sebagai kalor; sebuah belahan
memutus lilitan tertutupnya.

\textbf{23.} Terbuka: hanya $h$ yang meredam, $Q = 3.2$, ia berdenting. Terhubung singkat: ggl
baliknya mengalirkan arus lewat $R$, peredaman $B^2\ell^2/R$ ditambahkan,
$Q = 0.77$: ia berhenti seketika. (Cobalah.)

\textbf{24.} Sepanjang satu periode, energi tersimpannya kembali ke nilainya:
$\langle ui\rangle = \langle Ri^2\rangle + \langle hv^2\rangle$ --- masukan listrik $=$
kalor Joule dalam kumparannya $+$ rugi mekanisnya (bunyi yang terpancar dan
gesekan gantungannya); suku koplingnya $B\ell iv$ muncul dalam kedua
persamaannya dengan tanda berlawanan lalu saling meniadakan.

\textbf{25.} Dinamonya: Faraday pada kumparan yang berputar --- \qty{13.5}{V} pada
\qty{18}{km/h}, arusnya terbatasi pada \qty{1.35}{A} oleh induktansinya sendiri.
\omterm{prop:b1:induction:loudspeaker}{Pengeras suaranya}: \omterm{thm:b1:magnetostatics:laplace}{gaya Laplace} dan ggl balik --- \qty{3.1}{N.s/m} peredaman
listrik, \omterm{def:b1:heat-engines:efficiency}{efisiensi} $0.08\%$. Mikrofonnya: kumparan yang sama secara terbalik ---
\qty{5}{mV} tiap milimeter tiap sekon.
\end{solution}
